% HAND-WRITTEN fragment -- not generated by maestro2tex, but placed by it via the
% `order` list in configs/anatop_tb_tia.json.
%
% Every device parameter quoted here is a row of tab_tia_ssparams, extracted from the
% dcOpInfo of UnityFeedbackTB; every measured value is a row of tab_tia_amp. If the run
% is repeated, re-read both rather than editing numbers here.

\subsubsection{Small-signal model}\label{sss:tia-model}

The amplifier is the textbook two-stage Miller topology, and the textbook equations
describe it to within a few percent. Writing $g_{m1}$ for the input pair (M4),
$R_1 = (g_{o,\mathrm{M4}} + g_{o,\mathrm{M11}})^{-1}$ for the first-stage output node
$V_1$, and $g_{m2}$, $R_2 = (g_{o,\mathrm{M12}} + g_{o,\mathrm{M5}})^{-1}$ for the
second stage, the DC gain is the product of the two stage gains:

\begin{equation}\label{eq:tia-a0}
  A_0 = g_{m1} R_1 \cdot g_{m2} R_2 .
\end{equation}

With the extracted values of Table~\ref{tab:tia-ssparams} this is
$66.4 \times 249.3 = \num{16550}$, or \SI{84.38}{\decibel}, against
\SI{84.32}{\decibel} measured --- and the two agree to better than \SI{0.1}{\decibel}
at all five corners.

$C_c$ bridges the second stage, so it appears at $V_1$ multiplied by that stage's gain.
The resulting \SI{250}{\pico\farad} of effective capacitance on a \SI{2.1}{\mega\ohm}
node is the dominant pole, and it is the whole of the compensation:

\begin{equation}\label{eq:tia-poles}
  f_{p1} = \frac{1}{2\pi R_1 (1 + g_{m2} R_2) C_c},
  \qquad
  f_u = A_0 f_{p1} = \frac{g_{m1}}{2\pi C_c},
  \qquad
  f_{p2} \approx \frac{g_{m2}}{2\pi C_L}.
\end{equation}

These give \SI{301}{\hertz}, \SI{5.00}{\mega\hertz} and \SI{10.8}{\mega\hertz} against a
measured $f_{p1}$ of \SI{290}{\hertz} and a measured crossover of
\SI{6.39}{\mega\hertz}. Note the direction of the $f_u$ discrepancy: the loop crosses
\emph{above} the gain--bandwidth product, which a two-pole roll-off cannot do. A zero
does, and there is one.

\paragraph{The zero, and why $R_z$ exists.}
The $C_c$--$R_z$ branch is a feedforward path as well as a feedback path, and it puts a
zero in the numerator at

\begin{equation}\label{eq:tia-zero}
  f_z = \frac{1}{2\pi C_c \left( \dfrac{1}{g_{m2}} - R_z \right)} .
\end{equation}

Three regimes follow, and they are the reason the trim is there:

\begin{itemize}
\item $R_z = 0$: the zero is in the \emph{right} half plane at
$g_{m2}/2\pi C_c = \SI{53.8}{\mega\hertz}$. A RHP zero adds gain while subtracting
phase --- the classic failure mode of an uncompensated two-stage amplifier.
\item $R_z = 1/g_{m2} = \SI{2.96}{\kilo\ohm}$: the bracket vanishes and the zero moves
to infinity. The feedforward path is cancelled exactly.
\item $R_z > 1/g_{m2}$: the zero is in the \emph{left} half plane and adds phase lead.
Choosing
\begin{equation}\label{eq:tia-rz}
  R_z = \frac{1}{g_{m2}}\left(1 + \frac{C_L}{C_c}\right) = \SI{17.8}{\kilo\ohm}
\end{equation}
places it exactly on $f_{p2}$ and cancels the second pole.
\end{itemize}

This is what the 6-bit \texttt{LCBias} ladder adjusts. At the \texttt{LCBias}~=~15
setting simulated here the two largest ladder segments stay in circuit, giving an $R_z$
of order \SI{30}{\kilo\ohm} from the poly geometry --- above the cancellation value of
Equation~(\ref{eq:tia-rz}), so the zero lands \emph{below} $f_{p2}$. That is consistent
with both observations: the phase lead arrives before the second pole, which is why the
phase margin is \SI{88}{\degree} rather than the \SI{60}{\degree} a bare two-pole design
would give, and the extra gain pushes the crossover from \SI{5.0}{\mega\hertz} out to
\SI{6.4}{\mega\hertz}. ($R_z$ is inferred from the drawn geometry; it is not directly
observable in this run.) The phase that Figure~\ref{fig:tia-loopphase} plots is

\begin{equation}\label{eq:tia-phase}
  \angle A(jf) = -\arctan\frac{f}{f_{p1}} - \arctan\frac{f}{f_{p2}}
                 + \arctan\frac{f}{f_z},
  \qquad
  \mathrm{PM} = \ang{180} + \angle A(jf_u),
\end{equation}

the last term carrying the $+$ sign only because the zero is left-half-plane.

\paragraph{Large signal.}
The class-A second stage sets a hard ceiling on the rising edge: it can source only the
fixed current in M5, into $C_L$ and $C_c$ together,

\begin{equation}\label{eq:tia-slew}
  \mathrm{SR}^{+} \le \frac{I_{\mathrm{M5}}}{C_L + C_c}
  = \frac{\SI{14.4}{\micro\ampere}}{\SI{6.0}{\pico\farad}}
  = \SI{2.41}{\volt\per\micro\second}.
\end{equation}

The open-loop transient of Figure~\ref{fig:tia-slew} gives
\SI{2.20}{\volt\per\micro\second} and the closed-loop step of
Figure~\ref{fig:tia-step} slews at \SI{2.02}{\volt\per\micro\second}: the ceiling binds
in both. The falling edge has no such limit --- M12 is a driven device, not a current
source --- and is correspondingly faster at \SI{3.80}{\volt\per\micro\second}.
